\documentclass[10pt,a4paper]{amsart}
\usepackage{amsmath,amssymb,amsthm,graphicx}
\usepackage{verbatim}
\newenvironment{reference}{\comment}{\endcomment}
%\newenvironment{reference}{}{}
\newenvironment{slogan}{\comment}{\endcomment}
\newenvironment{history}{\comment}{\endcomment}

% For commutative diagrams we use Xy-pic
\usepackage[all]{xy}

% We use 2cell for 2-commutative diagrams.
\xyoption{2cell}
% Theorem environments.
%
\theoremstyle{plain}
\newtheorem{theorem}[subsection]{Theorem}
\newtheorem{proposition}[subsection]{Proposition}
\newtheorem{lemma}[subsection]{Lemma}

\theoremstyle{definition}
\newtheorem{definition}[subsection]{Definition}
\newtheorem{example}[subsection]{Example}
\newtheorem{exercise}[subsection]{Exercise}
\newtheorem{situation}[subsection]{Situation}

\theoremstyle{remark}
\newtheorem{remark}[subsection]{Remark}
\newtheorem{remarks}[subsection]{Remarks}

\numberwithin{equation}{subsection}

% Macros
%
\def\lim{\mathop{\mathrm{lim}}\nolimits}
\def\colim{\mathop{\mathrm{colim}}\nolimits}
\def\Spec{\mathop{\mathrm{Spec}}}
\def\Hom{\mathop{\mathrm{Hom}}\nolimits}
\def\Ext{\mathop{\mathrm{Ext}}\nolimits}
\def\SheafHom{\mathop{\mathcal{H}\!\mathit{om}}\nolimits}
\def\SheafExt{\mathop{\mathcal{E}\!\mathit{xt}}\nolimits}
\def\Sch{\mathit{Sch}}
\def\Mor{\mathop{\mathrm{Mor}}\nolimits}
\def\Ob{\mathop{\mathrm{Ob}}\nolimits}
\def\Sh{\mathop{\mathit{Sh}}\nolimits}
\def\NL{\mathop{N\!L}\nolimits}
\def\CH{\mathop{\mathrm{CH}}\nolimits}
\def\proetale{{pro\text{-}\acute{e}tale}}
\def\etale{{\acute{e}tale}}
\def\QCoh{\mathit{QCoh}}
\def\Ker{\mathop{\mathrm{Ker}}}
\def\Im{\mathop{\mathrm{Im}}}
\def\Coker{\mathop{\mathrm{Coker}}}
\def\Coim{\mathop{\mathrm{Coim}}}

% Boxtimes
%
\DeclareMathSymbol{\boxtimes}{\mathbin}{AMSa}{"02}

%
% Macros for moduli stacks/spaces
%
\def\QCohstack{\mathcal{QC}\!\mathit{oh}}
\def\Cohstack{\mathcal{C}\!\mathit{oh}}
\def\Spacesstack{\mathcal{S}\!\mathit{paces}}
\def\Quotfunctor{\mathrm{Quot}}
\def\Hilbfunctor{\mathrm{Hilb}}
\def\Curvesstack{\mathcal{C}\!\mathit{urves}}
\def\Polarizedstack{\mathcal{P}\!\mathit{olarized}}
\def\Complexesstack{\mathcal{C}\!\mathit{omplexes}}
% \Pic is the operator that assigns to X its picard group, usage \Pic(X)
% \Picardstack_{X/B} denotes the Picard stack of X over B
% \Picardfunctor_{X/B} denotes the Picard functor of X over B
\def\Pic{\mathop{\mathrm{Pic}}\nolimits}
\def\Picardstack{\mathcal{P}\!\mathit{ic}}
\def\Picardfunctor{\mathrm{Pic}}
\def\Deformationcategory{\mathcal{D}\!\mathit{ef}}
\title{Existence of transcendence bases and equality of their cardinalities}
\author{The Stacks Project Authors}
\date{stacks-project:fields.tex:local-snapshot-20261012; GNU FDL 1.2 or later}
\begin{document}
\maketitle
\noindent\textit{Editorial note.} All delivered mathematical text is the author's own, copied from the pinned Stacks Project chapter file fields.tex; no statement or proof step was rewritten, reordered or removed.\par\medskip
\noindent\textit{Editorial note.} The counted proof is the author's complete argument: The exchange step reuses the irreducible polynomial f: if g(beta\_1, alpha\_2, ..., alpha\_n) = 0 with g in F[x, y\_2, ..., y\_n] involved x, then beta\_1 would be algebraic over F(alpha\_2, ..., alpha\_n), contradicting the choice of f involving y\_1.\par\medskip
\noindent\textit{Editorial note.} The label lemma-transcendence-degree names the counted statement, and the proof environment immediately following it in the source is delivered verbatim as its proof.\par\medskip
\noindent\textit{Editorial note.} The counted statement is carried by the author's own lemma environment in the chapter source.\par\medskip
\noindent\textit{Allowed standard prerequisites (editorial).}\begin{itemize}
\item Context span 3585-3605: a transcendence basis of K/k is an algebraically independent collection \{x\_i\} such that K/k(x\_i; i in I) is algebraic
\item Zorn's lemma applied to the collection of algebraically independent subsets of G containing A; the union of a totally ordered family stays algebraically independent because polynomials involve only finitely many variables
\end{itemize}\medskip
\begin{definition}
\label{definition-transcendence}
Let $K/k$ be a field extension.
\begin{enumerate}
\item A collection of elements $\{x_i\}_{i \in I}$ of $K$ is called
{\it algebraically independent} over $k$ if the map
$$
k[X_i; i\in I] \longrightarrow K
$$
which maps $X_i$ to $x_i$ is injective.
\item The field of fractions of a polynomial ring
$k[x_i; i \in I]$ is denoted $k(x_i; i\in I)$.
\item A {\it purely transcendental extension} of $k$ is any
field extension $K/k$ isomorphic to the field of
fractions of a polynomial ring over $k$.
\item A {\it transcendence basis} of $K/k$ is a
collection of elements $\{x_i\}_{i \in I}$ which are
algebraically independent over $k$ and such that
the extension $K/k(x_i; i\in I)$ is algebraic.
\end{enumerate}
\end{definition}

\begin{lemma}
\label{lemma-transcendence-degree}
Let $E/F$ be a field extension. A transcendence basis of $E$ over $F$ exists.
Any two transcendence bases have the same cardinality.
\end{lemma}

\begin{proof}
Let $A$ be an algebraically independent subset of $E$. Let $G$ be a subset
of $E$ containing $A$ that generates $E/F$. We claim we can find a
transcendence basis $B$ such that $A \subset B \subset G$.
To prove this, consider the collection $\mathcal{B}$ of algebraically
independent subsets whose members are subsets of $G$ that contain $A$.
Define a partial ordering on $\mathcal{B}$ using inclusion.
Then $\mathcal{B}$ contains at least one element $A$.
The union of the elements of a totally ordered subset $T$ of $\mathcal{B}$
is an algebraically independent subset of $E$ over $F$ since any algebraic
dependence relation would have occurred in one of the elements of $T$
(since polynomials only involve finitely many variables). The union also
contains $A$ and is contained in $G$. By Zorn's lemma, there is a maximal
element $B \in \mathcal{B}$. Now we claim $E$ is algebraic over $F(B)$.
This is because if it wasn't then there would be an element
$f \in G$ transcendental over $F(B)$ since $F(G) = E$. Then
$B \cup\{f\}$ would be algebraically independent contradicting the
maximality of $B$. Thus $B$ is our transcendence basis.

\medskip\noindent
Let $B$ and $B'$ be two transcendence bases. Without loss of generality, we
can assume that $|B'| \leq |B|$. Now we divide the proof into two cases: the
first case is that $B$ is an infinite set. Then for each $\alpha \in B'$,
there is a finite set $B_{\alpha} \subset B$ such that $\alpha$ is algebraic
over
$F(B_{\alpha})$ since any algebraic dependence relation only uses finitely many
indeterminates. Then we define $B^* = \bigcup_{\alpha\in B'} B_{\alpha}$.
By construction, $B^* \subset B$, but we claim that in fact the two sets are
equal. To see this, suppose that they are not equal, say there is an element
$\beta \in B \setminus B^*$. We know $\beta$ is algebraic over $F(B')$ which
is algebraic over $F(B^*)$. Therefore $\beta$ is algebraic over $F(B^*)$, a
contradiction. So $|B| \leq |\bigcup_{\alpha \in B'} B_{\alpha}|$.
Now if $B'$ is finite, then so is $B$ so we can assume $B'$ is infinite;
this means
$$
|B| \leq |\bigcup\nolimits_{\alpha \in B'} B_{\alpha}| = |B'|
$$
because each $B_\alpha$ is finite and $B'$ is infinite. Therefore in the
infinite case, $|B| = |B'|$.

\medskip\noindent
Now we need to look at the case where $B$ is finite.
In this case, $B'$ is also finite, so suppose
$B = \{\alpha_1, \ldots, \alpha_n\}$ and
$B' = \{\beta_1, \ldots, \beta_m\}$ with $m \leq n$.
We perform induction on $m$: if $m = 0$ then $E/F$ is algebraic so
$B = \emptyset$ so $n = 0$. If $m > 0$, there is an irreducible polynomial
$f \in F[x, y_1, \ldots, y_n]$ such that
$f(\beta_1, \alpha_1, \ldots, \alpha_n) = 0$ and such that $x$ occurs in $f$.
Since $\beta_1$ is not algebraic over $F$, $f$ must involve some $y_i$
so without loss of generality, assume $f$ uses $y_1$.
Let $B^* = \{\beta_1, \alpha_2, \ldots, \alpha_n\}$.
We claim that $B^*$ is a basis for $E/F$. To prove this claim, we see that
we have a tower of algebraic extensions
$$
E/ F(B^*, \alpha_1) / F(B^*)
$$
since $\alpha_1$ is algebraic over $F(B^*)$.
Now we claim that $B^*$ (counting multiplicity of elements) is
algebraically independent over $F$ because if it weren't, then there would be an
irreducible $g\in F[x, y_2, \ldots, y_n]$ such that
$g(\beta_1, \alpha_2, \ldots, \alpha_n) = 0$
which must involve $x$ making $\beta_1$
algebraic over $F(\alpha_2, \ldots, \alpha_n)$ which would make $\alpha_1$
algebraic over $F(\alpha_2, \ldots, \alpha_n)$ which is impossible.
So this means that $\{\alpha_2, \ldots, \alpha_n\}$ and
$\{\beta_2, \ldots, \beta_m\}$ are bases for $E$ over $F(\beta_1)$
which means by induction, $m = n$.
\end{proof}

\end{document}
